\subsection{围绝经期性欲：变化规律与主动应对}

前面各节从"性欲低下的原因"角度提到围绝经期，这里把它单独展开——因为这不是一段"性欲注定走下坡"的岁月，而是一个\textbf{变化方式因人而异、且大有可为}的阶段。

\vspace{0.5em}
\noindent\textbf{激素变化如何影响性欲}：

\begin{itemize}
	\item \textbf{雌激素波动性下降}：是潮热、睡眠障碍、情绪波动的主因；而睡不好、情绪差，本身就直接压低性欲——很多女性误以为"我性冷淡了"，其实是\textbf{"我太累了"}
	\item \textbf{雄激素的缓坡下降}：与雌激素"断崖式"波动不同，女性的睾酮（参与性欲与性唤起）在围绝经期是平缓下降的，\textbf{不会一夜归零}
	\item \textbf{局部改变}：阴道黏膜变薄、润滑变慢、干涩与性交痛增多——"身体报警"让大脑把性与不适关联，性欲因此被条件反射式地抑制；这是\textbf{可逆的机械问题}，不是欲望消失
	\item \textbf{心理与社会因素}：容貌焦虑、"更年期等于衰老"的暗示、子女离家/照顾老人/婚姻倦怠叠加——这一阶段的心理负荷常被低估，其影响不亚于激素
\end{itemize}

\vspace{0.5em}
\noindent\textbf{个体差异远比"平均值"重要}：

\begin{itemize}
	\item 相当一部分女性报告围绝经期及绝经后\textbf{性欲回升、性满意度提高}——不再担心怀孕、不再被月经打扰、孩子长大有了二人世界，都是真实的"加项"
	\item 性欲下降明显、持续半年以上且带来痛苦者，才构成需要评估的问题（参考 HSDD 标准）；"比二十岁时少了"不是病，"想到性就烦躁回避"才值得就医
\end{itemize}

\vspace{0.5em}
\noindent\textbf{主动应对清单}：

\begin{itemize}
	\item \textbf{先解决"机械问题"}：水基/硅基润滑剂是第一选择；反复干涩、性交痛者可与医生讨论\textbf{局部（阴道）低剂量雌激素}——全身吸收极少，长期获益明确，是这一阶段的"性价比之王"
	\item \textbf{睡眠与潮热要管}：改善睡眠卫生、必要时评估激素治疗（MHT）——睡好了，欲望常常自己回来
	\item \textbf{延长前戏、放慢节奏}：围绝经期身体需要的唤起时间更长，这不是劣势，而是把"慢"变成新的亲密语言的机会
	\item \textbf{运动与盆底训练}：规律有氧与力量训练改善血流与情绪；凯格尔运动维持盆底肌力，改善润滑相关的舒适度与高潮质量
	\item \textbf{和伴侣一起更新"剧本"}：接纳身体的变化，坦率告诉对方"现在怎样更舒服"；把这一阶段当作\textbf{重新认识彼此}的机会，而不是告别
\end{itemize}

\begin{tcolorbox}[colback=pink!5!white,colframe=red!50!black,title=贴心小叮咛]
围绝经期的性欲波动是\textbf{激素、睡眠、情绪与关系}四件事的合成结果，而每一件都有办法应对。别用"年纪到了"一句话把自己打发掉——带着具体的困扰去妇科或更年期门诊，你会发现可选项比想象中多得多。
\end{tcolorbox}

